
\documentclass{article}
\usepackage{colortbl}
\usepackage{makecell}
\usepackage{multirow}
\usepackage{supertabular}

\begin{document}

\newcounter{utterance}

\centering \large Interaction Transcript for game `matchit\_ascii', experiment `similar\_grid\_2', episode 1 with glm{-}5.2.
\vspace{24pt}

{ \footnotesize  \setcounter{utterance}{1}
\setlength{\tabcolsep}{0pt}
\begin{supertabular}{c@{$\;$}|p{.15\linewidth}@{}p{.15\linewidth}p{.15\linewidth}p{.15\linewidth}p{.15\linewidth}p{.15\linewidth}}
    \# & $\;$A & \multicolumn{4}{c}{Game Master} & $\;\:$B\\
    \hline

    \theutterance \stepcounter{utterance}  
    & & \multicolumn{4}{p{0.6\linewidth}}{
        \cellcolor[rgb]{0.9,0.9,0.9}{
            \makecell[{{p{\linewidth}}}]{
                \texttt{\tiny{[P1$\langle$GM]}}
                \texttt{U neemt deel aan een gezamenlijk raadspel.} \\
\texttt{Het doel is om erachter te komen of dit raster en een andere afbeelding die alleen ik kan zien, hetzelfde zijn.} \\
\texttt{Beschrijf eerst uw raster. Daarna geef ik mijn beschrijving en kunnen we elkaar vragen stellen over de rasters om erachter te komen of ze hetzelfde zijn. Begin nu uw korte rasterbeschrijving met "BESCHRIJVING:" gevolgd door de beschrijving. Voeg niets anders toe. Dit is uw raster:} \\
\texttt{▢ ▢ X ▢ ▢} \\
\texttt{▢ ▢ X ▢ ▢} \\
\texttt{▢ ▢ X ▢ ▢} \\
\texttt{▢ ▢ X ▢ ▢} \\
\texttt{X X X X X} \\
            }
        }
    }
    & \\ \\

    \theutterance \stepcounter{utterance}  
    & \multicolumn{4}{p{0.6\linewidth}}{
        \cellcolor[rgb]{0.9,0.9,0.9}{
            \makecell[{{p{\linewidth}}}]{
                \texttt{\tiny{[P1$\rangle$GM]}}
                \texttt{BESCHRIJVING: Een 5x5 raster met een verticale lijn van vier 'X'en in de derde kolom op de eerste vier rijen, en een volledige horizontale rij van vijf 'X'en op de vijfde (onderste) rij. Alle andere vakjes zijn leeg ('▢').} \\
            }
        }
    }
    & & \\ \\

    \theutterance \stepcounter{utterance}  
    & & & \multicolumn{2}{p{0.3\linewidth}}{
        \cellcolor[rgb]{0.9,0.9,0.9}{
            \makecell[{{p{\linewidth}}}]{
                \texttt{\tiny{[GM$|$GM]}}
                \texttt{continue} \\
            }
        }
    }
    & & \\ \\

    \theutterance \stepcounter{utterance}  
    & & \multicolumn{4}{p{0.6\linewidth}}{
        \cellcolor[rgb]{0.9,0.9,0.9}{
            \makecell[{{p{\linewidth}}}]{
                \texttt{\tiny{[GM$\rangle$P2]}}
                \texttt{U neemt deel aan een gezamenlijk raadspel.} \\
\texttt{Het doel is om erachter te komen of dit raster en een andere afbeelding die alleen ik kan zien, hetzelfde zijn.} \\
\texttt{Beschrijf eerst uw raster. Daarna geef ik mijn beschrijving en kunnen we elkaar vragen stellen over de rasters om erachter te komen of ze hetzelfde zijn. Begin nu uw korte rasterbeschrijving met "BESCHRIJVING:" gevolgd door de beschrijving. Voeg niets anders toe. Dit is uw raster:} \\
\texttt{▢ ▢ X ▢ ▢} \\
\texttt{▢ ▢ X ▢ ▢} \\
\texttt{▢ ▢ ▢ ▢ ▢} \\
\texttt{▢ ▢ X X ▢} \\
\texttt{X X X X X} \\
            }
        }
    }
    & \\ \\

    \theutterance \stepcounter{utterance}  
    & & & \multicolumn{4}{p{0.6\linewidth}}{
        \cellcolor[rgb]{0.9,0.9,0.9}{
            \makecell[{{p{\linewidth}}}]{
                \texttt{\tiny{[GM$\langle$P2]}}
                \texttt{BESCHRIJVING: Het raster is 5x5 met een vorm gemaakt van X'en. Er is een verticale lijn van drie X'en in de derde kolom (rijen 1, 2 en 4). Op de vierde rij staan twee X'en naast elkaar in kolom 3 en 4. De vijfde rij is helemaal gevuld met X'en.} \\
            }
        }
    }
     \\ \\

    \theutterance \stepcounter{utterance}  
    & & & \multicolumn{2}{p{0.3\linewidth}}{
        \cellcolor[rgb]{0.9,0.9,0.9}{
            \makecell[{{p{\linewidth}}}]{
                \texttt{\tiny{[GM$|$GM]}}
                \texttt{continue} \\
            }
        }
    }
    & & \\ \\

    \theutterance \stepcounter{utterance}  
    & & \multicolumn{4}{p{0.6\linewidth}}{
        \cellcolor[rgb]{0.9,0.9,0.9}{
            \makecell[{{p{\linewidth}}}]{
                \texttt{\tiny{[GM$\rangle$P2]}}
                \texttt{Dit is mijn BESCHRIJVINGBESCHRIJVING: Een 5x5 raster met een verticale lijn van vier 'X'en in de derde kolom op de eerste vier rijen, en een volledige horizontale rij van vijf 'X'en op de vijfde (onderste) rij. Alle andere vakjes zijn leeg ('▢').} \\
\texttt{Stel nu een vraag om nieuwe aspecten van mijn raster te ontdekken die anders kunnen zijn dan jouw raster. Begin met "VRAAG:" en voeg niets anders toe.} \\
            }
        }
    }
    & \\ \\

    \theutterance \stepcounter{utterance}  
    & & & \multicolumn{4}{p{0.6\linewidth}}{
        \cellcolor[rgb]{0.9,0.9,0.9}{
            \makecell[{{p{\linewidth}}}]{
                \texttt{\tiny{[GM$\langle$P2]}}
                \texttt{VRAAG: Klopt het dat in jouw raster de vierde rij een X heeft in de derde kolom, maar niet in de vierde kolom?} \\
            }
        }
    }
     \\ \\

    \theutterance \stepcounter{utterance}  
    & & & \multicolumn{2}{p{0.3\linewidth}}{
        \cellcolor[rgb]{0.9,0.9,0.9}{
            \makecell[{{p{\linewidth}}}]{
                \texttt{\tiny{[GM$|$GM]}}
                \texttt{continue} \\
            }
        }
    }
    & & \\ \\

    \theutterance \stepcounter{utterance}  
    & & \multicolumn{4}{p{0.6\linewidth}}{
        \cellcolor[rgb]{0.9,0.9,0.9}{
            \makecell[{{p{\linewidth}}}]{
                \texttt{\tiny{[P1$\langle$GM]}}
                \texttt{Dit is mijn BESCHRIJVINGBESCHRIJVING: Het raster is 5x5 met een vorm gemaakt van X'en. Er is een verticale lijn van drie X'en in de derde kolom (rijen 1, 2 en 4). Op de vierde rij staan twee X'en naast elkaar in kolom 3 en 4. De vijfde rij is helemaal gevuld met X'en.} \\
\texttt{VRAAG: Klopt het dat in jouw raster de vierde rij een X heeft in de derde kolom, maar niet in de vierde kolom?Begin je antwoord met "ANTWOORD:" en voeg niets anders toe.} \\
            }
        }
    }
    & \\ \\

    \theutterance \stepcounter{utterance}  
    & \multicolumn{4}{p{0.6\linewidth}}{
        \cellcolor[rgb]{0.9,0.9,0.9}{
            \makecell[{{p{\linewidth}}}]{
                \texttt{\tiny{[P1$\rangle$GM]}}
                \texttt{ANTWOORD: Ja, in mijn raster heeft de vierde rij een X in de derde kolom en geen X in de vierde kolom.} \\
            }
        }
    }
    & & \\ \\

    \theutterance \stepcounter{utterance}  
    & & & \multicolumn{2}{p{0.3\linewidth}}{
        \cellcolor[rgb]{0.9,0.9,0.9}{
            \makecell[{{p{\linewidth}}}]{
                \texttt{\tiny{[GM$|$GM]}}
                \texttt{continue} \\
            }
        }
    }
    & & \\ \\

    \theutterance \stepcounter{utterance}  
    & & \multicolumn{4}{p{0.6\linewidth}}{
        \cellcolor[rgb]{0.9,0.9,0.9}{
            \makecell[{{p{\linewidth}}}]{
                \texttt{\tiny{[P1$\langle$GM]}}
                \texttt{Stel nu een vraag om nieuwe aspecten van mijn raster te ontdekken die anders kunnen zijn dan jouw raster. Begin met "VRAAG:" en voeg niets anders toe.} \\
            }
        }
    }
    & \\ \\

    \theutterance \stepcounter{utterance}  
    & \multicolumn{4}{p{0.6\linewidth}}{
        \cellcolor[rgb]{0.9,0.9,0.9}{
            \makecell[{{p{\linewidth}}}]{
                \texttt{\tiny{[P1$\rangle$GM]}}
                \texttt{VRAAG: Heeft jouw raster een X in de linkerbovenhoek (rij 1, kolom 1) of in de rechterbovenhoek (rij 1, kolom 5)?} \\
            }
        }
    }
    & & \\ \\

    \theutterance \stepcounter{utterance}  
    & & & \multicolumn{2}{p{0.3\linewidth}}{
        \cellcolor[rgb]{0.9,0.9,0.9}{
            \makecell[{{p{\linewidth}}}]{
                \texttt{\tiny{[GM$|$GM]}}
                \texttt{continue} \\
            }
        }
    }
    & & \\ \\

    \theutterance \stepcounter{utterance}  
    & & \multicolumn{4}{p{0.6\linewidth}}{
        \cellcolor[rgb]{0.9,0.9,0.9}{
            \makecell[{{p{\linewidth}}}]{
                \texttt{\tiny{[GM$\rangle$P2]}}
                \texttt{ANTWOORD: Ja, in mijn raster heeft de vierde rij een X in de derde kolom en geen X in de vierde kolom.} \\
\texttt{VRAAG: Heeft jouw raster een X in de linkerbovenhoek (rij 1, kolom 1) of in de rechterbovenhoek (rij 1, kolom 5)?Begin je antwoord met "ANTWOORD:" en voeg niets anders toe.} \\
            }
        }
    }
    & \\ \\

    \theutterance \stepcounter{utterance}  
    & & & \multicolumn{4}{p{0.6\linewidth}}{
        \cellcolor[rgb]{0.9,0.9,0.9}{
            \makecell[{{p{\linewidth}}}]{
                \texttt{\tiny{[GM$\langle$P2]}}
                \texttt{ANTWOORD: Nee, mijn raster heeft geen X in de linkerbovenhoek (rij 1, kolom 1) en ook geen X in de rechterbovenhoek (rij 1, kolom 5); beide vakjes zijn leeg (▢).} \\
            }
        }
    }
     \\ \\

    \theutterance \stepcounter{utterance}  
    & & & \multicolumn{2}{p{0.3\linewidth}}{
        \cellcolor[rgb]{0.9,0.9,0.9}{
            \makecell[{{p{\linewidth}}}]{
                \texttt{\tiny{[GM$|$GM]}}
                \texttt{continue} \\
            }
        }
    }
    & & \\ \\

    \theutterance \stepcounter{utterance}  
    & & \multicolumn{4}{p{0.6\linewidth}}{
        \cellcolor[rgb]{0.9,0.9,0.9}{
            \makecell[{{p{\linewidth}}}]{
                \texttt{\tiny{[GM$\rangle$P2]}}
                \texttt{Stel nu een vraag om nieuwe aspecten van mijn raster te ontdekken die anders kunnen zijn dan jouw raster. Begin met "VRAAG:" en voeg niets anders toe.} \\
            }
        }
    }
    & \\ \\

    \theutterance \stepcounter{utterance}  
    & & & \multicolumn{4}{p{0.6\linewidth}}{
        \cellcolor[rgb]{0.9,0.9,0.9}{
            \makecell[{{p{\linewidth}}}]{
                \texttt{\tiny{[GM$\langle$P2]}}
                \texttt{VRAAG: Zijn er in jouw raster op de vierde rij, behalve de X in de derde kolom, nog andere vakjes gevuld met een X?} \\
            }
        }
    }
     \\ \\

    \theutterance \stepcounter{utterance}  
    & & & \multicolumn{2}{p{0.3\linewidth}}{
        \cellcolor[rgb]{0.9,0.9,0.9}{
            \makecell[{{p{\linewidth}}}]{
                \texttt{\tiny{[GM$|$GM]}}
                \texttt{continue} \\
            }
        }
    }
    & & \\ \\

    \theutterance \stepcounter{utterance}  
    & & \multicolumn{4}{p{0.6\linewidth}}{
        \cellcolor[rgb]{0.9,0.9,0.9}{
            \makecell[{{p{\linewidth}}}]{
                \texttt{\tiny{[P1$\langle$GM]}}
                \texttt{ANTWOORD: Nee, mijn raster heeft geen X in de linkerbovenhoek (rij 1, kolom 1) en ook geen X in de rechterbovenhoek (rij 1, kolom 5); beide vakjes zijn leeg (▢).} \\
\texttt{VRAAG: Zijn er in jouw raster op de vierde rij, behalve de X in de derde kolom, nog andere vakjes gevuld met een X?Begin je antwoord met "ANTWOORD:" en voeg niets anders toe.} \\
            }
        }
    }
    & \\ \\

    \theutterance \stepcounter{utterance}  
    & \multicolumn{4}{p{0.6\linewidth}}{
        \cellcolor[rgb]{0.9,0.9,0.9}{
            \makecell[{{p{\linewidth}}}]{
                \texttt{\tiny{[P1$\rangle$GM]}}
                \texttt{ANTWOORD: Nee, op de vierde rij is in mijn raster alleen het derde vakje gevuld met een X; de andere vakjes op die rij zijn leeg (▢).} \\
            }
        }
    }
    & & \\ \\

    \theutterance \stepcounter{utterance}  
    & & & \multicolumn{2}{p{0.3\linewidth}}{
        \cellcolor[rgb]{0.9,0.9,0.9}{
            \makecell[{{p{\linewidth}}}]{
                \texttt{\tiny{[GM$|$GM]}}
                \texttt{continue} \\
            }
        }
    }
    & & \\ \\

    \theutterance \stepcounter{utterance}  
    & & \multicolumn{4}{p{0.6\linewidth}}{
        \cellcolor[rgb]{0.9,0.9,0.9}{
            \makecell[{{p{\linewidth}}}]{
                \texttt{\tiny{[P1$\langle$GM]}}
                \texttt{Stel nu een vraag om nieuwe aspecten van mijn raster te ontdekken die anders kunnen zijn dan jouw raster. Begin met "VRAAG:" en voeg niets anders toe.} \\
            }
        }
    }
    & \\ \\

    \theutterance \stepcounter{utterance}  
    & \multicolumn{4}{p{0.6\linewidth}}{
        \cellcolor[rgb]{0.9,0.9,0.9}{
            \makecell[{{p{\linewidth}}}]{
                \texttt{\tiny{[P1$\rangle$GM]}}
                \texttt{VRAAG: In jouw beschrijving staat dat de verticale lijn in de derde kolom op rijen 1, 2 en 4 ligt. Betekent dit dat rij 3 in de derde kolom leeg (▢) is, en zijn er verder op rij 3 nog andere X'en te vinden?} \\
            }
        }
    }
    & & \\ \\

    \theutterance \stepcounter{utterance}  
    & & & \multicolumn{2}{p{0.3\linewidth}}{
        \cellcolor[rgb]{0.9,0.9,0.9}{
            \makecell[{{p{\linewidth}}}]{
                \texttt{\tiny{[GM$|$GM]}}
                \texttt{continue} \\
            }
        }
    }
    & & \\ \\

    \theutterance \stepcounter{utterance}  
    & & \multicolumn{4}{p{0.6\linewidth}}{
        \cellcolor[rgb]{0.9,0.9,0.9}{
            \makecell[{{p{\linewidth}}}]{
                \texttt{\tiny{[GM$\rangle$P2]}}
                \texttt{ANTWOORD: Nee, op de vierde rij is in mijn raster alleen het derde vakje gevuld met een X; de andere vakjes op die rij zijn leeg (▢).} \\
\texttt{VRAAG: In jouw beschrijving staat dat de verticale lijn in de derde kolom op rijen 1, 2 en 4 ligt. Betekent dit dat rij 3 in de derde kolom leeg (▢) is, en zijn er verder op rij 3 nog andere X'en te vinden?Begin je antwoord met "ANTWOORD:" en voeg niets anders toe.} \\
            }
        }
    }
    & \\ \\

    \theutterance \stepcounter{utterance}  
    & & & \multicolumn{4}{p{0.6\linewidth}}{
        \cellcolor[rgb]{0.9,0.9,0.9}{
            \makecell[{{p{\linewidth}}}]{
                \texttt{\tiny{[GM$\langle$P2]}}
                \texttt{ANTWOORD: Ja, rij 3 in de derde kolom is leeg (▢) en er zijn verder op rij 3 helemaal geen X'en; de hele rij is leeg.} \\
            }
        }
    }
     \\ \\

    \theutterance \stepcounter{utterance}  
    & & & \multicolumn{2}{p{0.3\linewidth}}{
        \cellcolor[rgb]{0.9,0.9,0.9}{
            \makecell[{{p{\linewidth}}}]{
                \texttt{\tiny{[GM$|$GM]}}
                \texttt{continue} \\
            }
        }
    }
    & & \\ \\

    \theutterance \stepcounter{utterance}  
    & & \multicolumn{4}{p{0.6\linewidth}}{
        \cellcolor[rgb]{0.9,0.9,0.9}{
            \makecell[{{p{\linewidth}}}]{
                \texttt{\tiny{[GM$\rangle$P2]}}
                \texttt{Stel nu een vraag om nieuwe aspecten van mijn raster te ontdekken die anders kunnen zijn dan jouw raster. Begin met "VRAAG:" en voeg niets anders toe.} \\
            }
        }
    }
    & \\ \\

    \theutterance \stepcounter{utterance}  
    & & & \multicolumn{4}{p{0.6\linewidth}}{
        \cellcolor[rgb]{0.9,0.9,0.9}{
            \makecell[{{p{\linewidth}}}]{
                \texttt{\tiny{[GM$\langle$P2]}}
                \texttt{VRAAG: Zijn er in jouw raster op de eerste en tweede rij nog andere X'en aanwezig, behalve de X in de derde kolom?} \\
            }
        }
    }
     \\ \\

    \theutterance \stepcounter{utterance}  
    & & & \multicolumn{2}{p{0.3\linewidth}}{
        \cellcolor[rgb]{0.9,0.9,0.9}{
            \makecell[{{p{\linewidth}}}]{
                \texttt{\tiny{[GM$|$GM]}}
                \texttt{continue} \\
            }
        }
    }
    & & \\ \\

    \theutterance \stepcounter{utterance}  
    & & \multicolumn{4}{p{0.6\linewidth}}{
        \cellcolor[rgb]{0.9,0.9,0.9}{
            \makecell[{{p{\linewidth}}}]{
                \texttt{\tiny{[P1$\langle$GM]}}
                \texttt{ANTWOORD: Ja, rij 3 in de derde kolom is leeg (▢) en er zijn verder op rij 3 helemaal geen X'en; de hele rij is leeg.} \\
\texttt{VRAAG: Zijn er in jouw raster op de eerste en tweede rij nog andere X'en aanwezig, behalve de X in de derde kolom?Begin je antwoord met "ANTWOORD:" en voeg niets anders toe.} \\
            }
        }
    }
    & \\ \\

    \theutterance \stepcounter{utterance}  
    & \multicolumn{4}{p{0.6\linewidth}}{
        \cellcolor[rgb]{0.9,0.9,0.9}{
            \makecell[{{p{\linewidth}}}]{
                \texttt{\tiny{[P1$\rangle$GM]}}
                \texttt{ANTWOORD: Nee, op de eerste en tweede rij is in mijn raster alleen het derde vakje gevuld met een X; alle andere vakjes op die rijen zijn leeg (▢).} \\
            }
        }
    }
    & & \\ \\

    \theutterance \stepcounter{utterance}  
    & & & \multicolumn{2}{p{0.3\linewidth}}{
        \cellcolor[rgb]{0.9,0.9,0.9}{
            \makecell[{{p{\linewidth}}}]{
                \texttt{\tiny{[GM$|$GM]}}
                \texttt{continue} \\
            }
        }
    }
    & & \\ \\

    \theutterance \stepcounter{utterance}  
    & & \multicolumn{4}{p{0.6\linewidth}}{
        \cellcolor[rgb]{0.9,0.9,0.9}{
            \makecell[{{p{\linewidth}}}]{
                \texttt{\tiny{[P1$\langle$GM]}}
                \texttt{Stel nu een vraag om nieuwe aspecten van mijn raster te ontdekken die anders kunnen zijn dan jouw raster. Begin met "VRAAG:" en voeg niets anders toe.} \\
            }
        }
    }
    & \\ \\

    \theutterance \stepcounter{utterance}  
    & \multicolumn{4}{p{0.6\linewidth}}{
        \cellcolor[rgb]{0.9,0.9,0.9}{
            \makecell[{{p{\linewidth}}}]{
                \texttt{\tiny{[P1$\rangle$GM]}}
                \texttt{VRAAG: Je noemde eerder een verticale lijn op "rijen 1, 2 en 4" en een X op "rij 4" in kolom 3. Zijn de vakjes in de derde kolom op rij 1 en rij 2 in jouw raster gevuld met een X, of zijn die leeg?} \\
            }
        }
    }
    & & \\ \\

    \theutterance \stepcounter{utterance}  
    & & & \multicolumn{2}{p{0.3\linewidth}}{
        \cellcolor[rgb]{0.9,0.9,0.9}{
            \makecell[{{p{\linewidth}}}]{
                \texttt{\tiny{[GM$|$GM]}}
                \texttt{continue} \\
            }
        }
    }
    & & \\ \\

    \theutterance \stepcounter{utterance}  
    & & \multicolumn{4}{p{0.6\linewidth}}{
        \cellcolor[rgb]{0.9,0.9,0.9}{
            \makecell[{{p{\linewidth}}}]{
                \texttt{\tiny{[GM$\rangle$P2]}}
                \texttt{ANTWOORD: Nee, op de eerste en tweede rij is in mijn raster alleen het derde vakje gevuld met een X; alle andere vakjes op die rijen zijn leeg (▢).} \\
\texttt{VRAAG: Je noemde eerder een verticale lijn op "rijen 1, 2 en 4" en een X op "rij 4" in kolom 3. Zijn de vakjes in de derde kolom op rij 1 en rij 2 in jouw raster gevuld met een X, of zijn die leeg?Begin je antwoord met "ANTWOORD:" en voeg niets anders toe.} \\
            }
        }
    }
    & \\ \\

    \theutterance \stepcounter{utterance}  
    & & & \multicolumn{4}{p{0.6\linewidth}}{
        \cellcolor[rgb]{0.9,0.9,0.9}{
            \makecell[{{p{\linewidth}}}]{
                \texttt{\tiny{[GM$\langle$P2]}}
                \texttt{ANTWOORD: De vakjes in de derde kolom op rij 1 en rij 2 in mijn raster zijn gevuld met een X.} \\
            }
        }
    }
     \\ \\

    \theutterance \stepcounter{utterance}  
    & & & \multicolumn{2}{p{0.3\linewidth}}{
        \cellcolor[rgb]{0.9,0.9,0.9}{
            \makecell[{{p{\linewidth}}}]{
                \texttt{\tiny{[GM$|$GM]}}
                \texttt{continue} \\
            }
        }
    }
    & & \\ \\

    \theutterance \stepcounter{utterance}  
    & & \multicolumn{4}{p{0.6\linewidth}}{
        \cellcolor[rgb]{0.9,0.9,0.9}{
            \makecell[{{p{\linewidth}}}]{
                \texttt{\tiny{[GM$\rangle$P2]}}
                \texttt{Neem nu een beslissing. Wat denk je: zijn jouw raster en het andere beschreven raster hetzelfde? Schrijf "BESLISSING: hetzelfde raster" als je denkt dat het hetzelfde raster is of "BESLISSING: verschillende rasters" als je denkt dat het verschillende rasters zijn. Voeg niets anders toe.} \\
            }
        }
    }
    & \\ \\

    \theutterance \stepcounter{utterance}  
    & & & \multicolumn{4}{p{0.6\linewidth}}{
        \cellcolor[rgb]{0.9,0.9,0.9}{
            \makecell[{{p{\linewidth}}}]{
                \texttt{\tiny{[GM$\langle$P2]}}
                \texttt{BESLISSING: verschillende rasters} \\
            }
        }
    }
     \\ \\

    \theutterance \stepcounter{utterance}  
    & & & \multicolumn{2}{p{0.3\linewidth}}{
        \cellcolor[rgb]{0.9,0.9,0.9}{
            \makecell[{{p{\linewidth}}}]{
                \texttt{\tiny{[GM$|$GM]}}
                \texttt{continue} \\
            }
        }
    }
    & & \\ \\

    \theutterance \stepcounter{utterance}  
    & & & \multicolumn{2}{p{0.3\linewidth}}{
        \cellcolor[rgb]{0.9,0.9,0.9}{
            \makecell[{{p{\linewidth}}}]{
                \texttt{\tiny{[GM$|$GM]}}
                \texttt{success} \\
            }
        }
    }
    & & \\ \\

    \theutterance \stepcounter{utterance}  
    & & \multicolumn{4}{p{0.6\linewidth}}{
        \cellcolor[rgb]{0.9,0.9,0.9}{
            \makecell[{{p{\linewidth}}}]{
                \texttt{\tiny{[P1$\langle$GM]}}
                \texttt{ANTWOORD: De vakjes in de derde kolom op rij 1 en rij 2 in mijn raster zijn gevuld met een X.} \\
\texttt{Neem nu een beslissing. Wat denk je: zijn jouw raster en het andere beschreven raster hetzelfde? Schrijf "BESLISSING: hetzelfde raster" als je denkt dat het hetzelfde raster is of "BESLISSING: verschillende rasters" als je denkt dat het verschillende rasters zijn. Voeg niets anders toe.} \\
            }
        }
    }
    & \\ \\

    \theutterance \stepcounter{utterance}  
    & \multicolumn{4}{p{0.6\linewidth}}{
        \cellcolor[rgb]{0.9,0.9,0.9}{
            \makecell[{{p{\linewidth}}}]{
                \texttt{\tiny{[P1$\rangle$GM]}}
                \texttt{BESLISSING: verschillende rasters} \\
            }
        }
    }
    & & \\ \\

    \theutterance \stepcounter{utterance}  
    & & & \multicolumn{2}{p{0.3\linewidth}}{
        \cellcolor[rgb]{0.9,0.9,0.9}{
            \makecell[{{p{\linewidth}}}]{
                \texttt{\tiny{[GM$|$GM]}}
                \texttt{continue} \\
            }
        }
    }
    & & \\ \\

    \theutterance \stepcounter{utterance}  
    & & & \multicolumn{2}{p{0.3\linewidth}}{
        \cellcolor[rgb]{0.9,0.9,0.9}{
            \makecell[{{p{\linewidth}}}]{
                \texttt{\tiny{[GM$|$GM]}}
                \texttt{success} \\
            }
        }
    }
    & & \\ \\

\end{supertabular}
}

\end{document}
